% GENERATED FILE. Edit canonical lessons, then run npm run generate:books.
% canonical-source-hash: fnv1a64:c88d444eb6f57760
% canonical-lessons: KA-C224-hoge, KA-C224-dhulu, KA-C224-pravaha, KA-C224-prakrti, KA-C224-prani, KA-R224-first-pass-words-from-chapters-216-219, KA-R224-second-pass-words-from-chapters-220-224

\chapter{Smoke, Dust, Flood, Nature, Animal}
\label{ch:ka-smoke-dust-flood-nature-animal}
\hlchaptermodality{224}

\begin{chapteropening}
\textbf{By the end of this chapter:} \emph{I can say smoke, dust, a flood, nature and an animal.}

Complete the last lesson of chapter 224: I can say smoke, dust, a flood, nature and an animal.
\end{chapteropening}

\section[\texorpdfstring{\kn{ಹೊಗೆ} (hoge)}{ಹೊಗೆ (hoge)}]{\kn{ಹೊಗೆ} (hoge) — smoke}

\label{lesson:KA-C224-hoge}



\begin{quote}
Before the new one: say the Kannada for an identity card, then the Kannada for a cheque.
\end{quote}

\subsection*{You'll want to know: \kn{ಹೊಗೆ}}
\textbf{\kn{ಹೊಗೆ}} — \emph{hoge} — \textquotedblleft{}smoke\textquotedblright{}.

\begin{grammarlens}[title={What you've built}]
One more word for this chapter.
\end{grammarlens}

\subsection*{Your turn}
\begin{itemize}
\raggedright
  \item \emph{Say it:} \emph{hoge}
  \item \emph{Say it:} \emph{hoge}, once more
  \item \emph{Say it:} say \emph{hoge}
  \item \emph{Recall:} say the Kannada for an identity card, then the Kannada for a cheque, then say \emph{hoge} again
\end{itemize}

\begin{tcolorbox}[breakable,colback=teal!4,colframe=teal!35!black,title={Before you move on}]
What does \kn{ಹೊಗೆ} mean? (\textquotedblleft{}Smoke\textquotedblright{}.) Say it once more.
\end{tcolorbox}
\section[\texorpdfstring{\kn{ಧೂಳು} (dhūḷu)}{ಧೂಳು (dhūḷu)}]{\kn{ಧೂಳು} (dhūḷu) — dust}

\label{lesson:KA-C224-dhulu}



\begin{quote}
Before the new one: say the Kannada for a cheque, then the Kannada for smoke.
\end{quote}

\subsection*{You'll want to know: \kn{ಧೂಳು}}
\textbf{\kn{ಧೂಳು}} — \emph{dhūḷu} — \textquotedblleft{}dust\textquotedblright{}.

\begin{grammarlens}[title={What you've built}]
One more word for this chapter.
\end{grammarlens}

\subsection*{Your turn}
\begin{itemize}
\raggedright
  \item \emph{Say it:} \emph{dhūḷu}
  \item \emph{Say it:} \emph{dhūḷu}, once more
  \item \emph{Say it:} say \emph{dhūḷu}
  \item \emph{Recall:} say the Kannada for a cheque, then the Kannada for smoke, then say \emph{dhūḷu} again
\end{itemize}

\begin{tcolorbox}[breakable,colback=teal!4,colframe=teal!35!black,title={Before you move on}]
What does \kn{ಧೂಳು} mean? (\textquotedblleft{}Dust\textquotedblright{}.) Say it once more.
\end{tcolorbox}
\section[\texorpdfstring{\kn{ಪ್ರವಾಹ} (pravāha)}{ಪ್ರವಾಹ (pravāha)}]{\kn{ಪ್ರವಾಹ} (pravāha) — a flood}

\label{lesson:KA-C224-pravaha}



\begin{quote}
Before the new one: say the Kannada for smoke, then the Kannada for dust.
\end{quote}

\subsection*{You'll want to know: \kn{ಪ್ರವಾಹ}}
\textbf{\kn{ಪ್ರವಾಹ}} — \emph{pravāha} — \textquotedblleft{}a flood\textquotedblright{}.

\begin{grammarlens}[title={What you've built}]
One more word for this chapter.
\end{grammarlens}

\subsection*{Your turn}
\begin{itemize}
\raggedright
  \item \emph{Say it:} \emph{pravāha}
  \item \emph{Say it:} \emph{pravāha}, once more
  \item \emph{Say it:} say \emph{pravāha}
  \item \emph{Recall:} say the Kannada for smoke, then the Kannada for dust, then say \emph{pravāha} again
\end{itemize}

\begin{tcolorbox}[breakable,colback=teal!4,colframe=teal!35!black,title={Before you move on}]
What does \kn{ಪ್ರವಾಹ} mean? (\textquotedblleft{}A flood\textquotedblright{}.) Say it once more.
\end{tcolorbox}
\section[\texorpdfstring{\kn{ಪ್ರಕೃತಿ} (prakṛti)}{ಪ್ರಕೃತಿ (prakṛti)}]{\kn{ಪ್ರಕೃತಿ} (prakṛti) — nature}

\label{lesson:KA-C224-prakrti}



\begin{quote}
Before the new one: say the Kannada for dust, then the Kannada for a flood.
\end{quote}

\subsection*{You'll want to know: \kn{ಪ್ರಕೃತಿ}}
\textbf{\kn{ಪ್ರಕೃತಿ}} — \emph{prakṛti} — \textquotedblleft{}nature\textquotedblright{}.

\begin{grammarlens}[title={What you've built}]
One more word for this chapter.
\end{grammarlens}

\subsection*{Your turn}
\begin{itemize}
\raggedright
  \item \emph{Say it:} \emph{prakṛti}
  \item \emph{Say it:} \emph{prakṛti}, once more
  \item \emph{Say it:} say \emph{prakṛti}
  \item \emph{Recall:} say the Kannada for dust, then the Kannada for a flood, then say \emph{prakṛti} again
\end{itemize}

\begin{tcolorbox}[breakable,colback=teal!4,colframe=teal!35!black,title={Before you move on}]
What does \kn{ಪ್ರಕೃತಿ} mean? (\textquotedblleft{}Nature\textquotedblright{}.) Say it once more.
\end{tcolorbox}
\section[\texorpdfstring{\kn{ಪ್ರಾಣಿ} (prāṇi)}{ಪ್ರಾಣಿ (prāṇi)}]{\kn{ಪ್ರಾಣಿ} (prāṇi) — an animal}

\label{lesson:KA-C224-prani}



\begin{quote}
Before the new one: say the Kannada for a flood, then the Kannada for nature.
\end{quote}

\subsection*{You'll want to know: \kn{ಪ್ರಾಣಿ}}
\textbf{\kn{ಪ್ರಾಣಿ}} — \emph{prāṇi} — \textquotedblleft{}an animal\textquotedblright{}.

\begin{grammarlens}[title={What you've built}]
One more word for this chapter.
\end{grammarlens}

\subsection*{Your turn}
\begin{itemize}
\raggedright
  \item \emph{Say it:} \emph{prāṇi}
  \item \emph{Say it:} \emph{prāṇi}, once more
  \item \emph{Say it:} say \emph{prāṇi}
  \item \emph{Recall:} say the Kannada for a flood, then the Kannada for nature, then say \emph{prāṇi} again
\end{itemize}

\begin{tcolorbox}[breakable,colback=teal!4,colframe=teal!35!black,title={Before you move on}]
What does \kn{ಪ್ರಾಣಿ} mean? (\textquotedblleft{}An animal\textquotedblright{}.) Say it once more.
\end{tcolorbox}
\section[(dialogue)]{Words from chapters 216-219 — a first pass}

\label{lesson:KA-R224-first-pass-words-from-chapters-216-219}



\begin{quote}
Say the words for nature and an animal.
\end{quote}

\subsection*{Your turn}
Say these aloud:

\begin{itemize}
\raggedright
  \item success, freedom, justice, responsibility, a duty
  \item cheap, empty, tight, loose, thick
  \item ready, available, clear, honest, generous
  \item ordinary, blunt, smooth, rough, sticky
\end{itemize}

\begin{tcolorbox}[breakable,colback=teal!4,colframe=teal!35!black,title={Before you move on}]
Success? (\textbf{\kn{ಯಶಸ್ಸು}}, \emph{yaśassu}.) Freedom? (\textbf{\kn{ಸ್ವಾತಂತ್ರ್ಯ}}, \emph{svātantrya}.) Justice? (\textbf{\kn{ನ್ಯಾಯ}}, \emph{nyāya}.) Responsibility? (\textbf{\kn{ಜವಾಬ್ದಾರಿ}}, \emph{javābdāri}.) A duty? (\textbf{\kn{ಕರ್ತವ್ಯ}}, \emph{kartavya}.) Cheap? (\textbf{\kn{ಅಗ್ಗ}}, \emph{agga}.) Empty? (\textbf{\kn{ಖಾಲಿ}}, \emph{khāli}.) Tight? (\textbf{\kn{ಬಿಗಿ}}, \emph{bigi}.) Loose? (\textbf{\kn{ಸಡಿಲ}}, \emph{saḍila}.) Thick? (\textbf{\kn{ದಪ್ಪ}}, \emph{dappa}.) Ready? (\textbf{\kn{ಸಿದ್ಧ}}, \emph{siddha}.) Available? (\textbf{\kn{ಲಭ್ಯ}}, \emph{labhya}.) Clear? (\textbf{\kn{ಸ್ಪಷ್ಟ}}, \emph{spaṣṭa}.) Honest? (\textbf{\kn{ಪ್ರಾಮಾಣಿಕ}}, \emph{prāmāṇika}.) Generous? (\textbf{\kn{ಉದಾರ}}, \emph{udāra}.) Ordinary? (\textbf{\kn{ಸಾಮಾನ್ಯ}}, \emph{sāmānya}.) Blunt? (\textbf{\kn{ಮೊಂಡು}}, \emph{moṇḍu}.) Smooth? (\textbf{\kn{ನುಣುಪು}}, \emph{nuṇupu}.) Rough? (\textbf{\kn{ಒರಟು}}, \emph{oraṭu}.) Sticky? (\textbf{\kn{ಜಿಗುಟು}}, \emph{jiguṭu}.)
\end{tcolorbox}
\section[(dialogue)]{Words from chapters 220-224 — a second pass}

\label{lesson:KA-R224-second-pass-words-from-chapters-220-224}



\begin{quote}
Say the words for a lid and a cupboard.
\end{quote}

\subsection*{Your turn}
Say these aloud:

\begin{itemize}
\raggedright
  \item a lid, a cupboard, a fridge, a fan, rubbish
  \item an apple, an orange, a papaya, a watermelon, a chapati
  \item trousers, a long skirt, a blouse, a vest, shoes
  \item a form, a certificate, a passport, an identity card, a cheque
  \item smoke, dust, a flood, nature, an animal
\end{itemize}

\begin{tcolorbox}[breakable,colback=teal!4,colframe=teal!35!black,title={Before you move on}]
A lid? (\textbf{\kn{ಮುಚ್ಚಳ}}, \emph{muccaḷa}.) A cupboard? (\textbf{\kn{ಬೀರು}}, \emph{bīru}.) A fridge? (\textbf{\kn{ಫ್ರಿಜ್}}, \emph{phrij}.) A fan? (\textbf{\kn{ಫ್ಯಾನ್}}, \emph{phyān}.) Rubbish? (\textbf{\kn{ಕಸ}}, \emph{kasa}.) An apple? (\textbf{\kn{ಸೇಬು}}, \emph{sēbu}.) An orange? (\textbf{\kn{ಕಿತ್ತಳೆ}}, \emph{kittaḷe}.) A papaya? (\textbf{\kn{ಪಪ್ಪಾಯಿ}}, \emph{pappāyi}.) A watermelon? (\textbf{\kn{ಕಲ್ಲಂಗಡಿ}}, \emph{kallaṅgaḍi}.) A chapati? (\textbf{\kn{ಚಪಾತಿ}}, \emph{capāti}.) Trousers? (\textbf{\kn{ಪ್ಯಾಂಟ್}}, \emph{pyāṇṭ}.) A long skirt? (\textbf{\kn{ಲಂಗ}}, \emph{laṅga}.) A blouse? (\textbf{\kn{ರವಿಕೆ}}, \emph{ravike}.) A vest? (\textbf{\kn{ಬನಿಯನ್}}, \emph{baniyan}.) Shoes? (\textbf{\kn{ಬೂಟು}}, \emph{būṭu}.) A form? (\textbf{\kn{ನಮೂನೆ}}, \emph{namūne}.) A certificate? (\textbf{\kn{ಪ್ರಮಾಣಪತ್ರ}}, \emph{pramāṇapatra}.) A passport? (\textbf{\kn{ಪಾಸ್ಪೋರ್ಟ್}}, \emph{pāspōrṭ}.) An identity card? (\textbf{\kn{ಗುರುತಿನ} \kn{ಚೀಟಿ}}, \emph{gurutina cīṭi}.) A cheque? (\textbf{\kn{ಚೆಕ್}}, \emph{cek}.) Smoke? (\textbf{\kn{ಹೊಗೆ}}, \emph{hoge}.) Dust? (\textbf{\kn{ಧೂಳು}}, \emph{dhūḷu}.) A flood? (\textbf{\kn{ಪ್ರವಾಹ}}, \emph{pravāha}.) Nature? (\textbf{\kn{ಪ್ರಕೃತಿ}}, \emph{prakṛti}.) An animal? (\textbf{\kn{ಪ್ರಾಣಿ}}, \emph{prāṇi}.)
\end{tcolorbox}
